\section{The elicitation protocol, and one record end to end}
\label{app:worked}

This appendix gives the exact prompts behind every record, then one complete record
from the Mistral-Small-24B $\times$ MATH cell, summarized first and then given in full.

\begin{center}
\footnotesize
\fbox{\begin{minipage}{0.96\linewidth}
\textbf{One record, Mistral-Small-24B $\times$ MATH.} \emph{What integer $n$ satisfies $0\le n<18$ and $n\equiv -11213141 \pmod{18}$?}
The model sets the problem up correctly and gets $18\times622953 = 11213154$ right, then writes $r = -11213141 + 11213154 = 3$ and answers $3$.
That subtraction gives $13$, which is the gold answer.
Stated confidence: $90$ before the attempt, $95$ after it.
Other readouts of the same attempt: P(True), the probability the model assigns to ``True'' when asked whether its answer is correct, $0.32$; mean answer log-probability $-0.05$.
\end{minipage}}
\end{center}

\noindent
A single arithmetic slip inside a fluent derivation. Having just made the error, the model
revises its stated confidence upward, while P(True) on the same transcript assigns only
0.32 probability to correctness. The rest of this appendix gives the prompts and the full transcript.

\subsection{Prompts}

Each record is four forward passes: the three turns of \cref{sec:reading} plus one read taken
beforehand, at the last prompt token of the bare question, before any elicitation. The confidence turns
run with thinking off so the read lands on a clean one-line reply; the attempt turn runs
with thinking on for hybrid reasoners. The fixed \texttt{Confidence: N} format is what makes
the read site content-neutral: the digit token has identical local context on every record,
so layer-0 activations carry no signal by construction.

\paragraph{System prompt (MATH).}
\begin{quote}\ttfamily\footnotesize
Solve the problem step by step. Put your final answer in \textbackslash boxed\{\}.
\end{quote}

\paragraph{Turn 1, prospective confidence.} The problem is followed by:
\begin{quote}\ttfamily\footnotesize
Before answering, estimate how likely you are to answer this question correctly. Reply with EXACTLY one line and nothing else:\\
Confidence: N\\
where N is an integer from 0 (certain to be wrong) to 100 (certain to be correct). Do NOT answer yet.
\end{quote}

\paragraph{Turn 2, the attempt.} The turn-1 reply is left in context as the assistant message, then:
\begin{quote}\ttfamily\footnotesize
Now give your final answer inside \textbackslash boxed\{\}.
\end{quote}

\paragraph{Turn 3, retrospective confidence.} The attempt is left in context, then:
\begin{quote}\ttfamily\footnotesize
Now assess the answer you just gave. On a scale of 0 to 100, how confident are you that it is correct? Reply with exactly one line and nothing else:\\
Confidence: N
\end{quote}

\paragraph{P(True) baseline.} A separate teacher-forced branch off turn 2, no generation.
Next-token mass on \texttt{True} versus \texttt{False}:
\begin{quote}\ttfamily\footnotesize
Consider your final answer above. Is it correct? Reply with exactly one word: True or False.
\end{quote}

Grading is external and is never shown to the model, so the turn-3 revision is blind
self-assessment.

\subsection{One record, verbatim}

Record \texttt{math\_test\_number\_theory\_753}, sample 2 of 5.

\paragraph{Problem.} \emph{What integer $n$ satisfies $0 \le n < 18$ and $n \equiv -11213141 \pmod{18}$?}

\paragraph{Turn 1 reply.} \texttt{Confidence: 90}

\paragraph{Turn 2 reply (complete).}
\begin{quote}\footnotesize
First, we need to find the remainder when $-11213141$ is divided by $18$.

1. Calculate $-11213141 \bmod 18$:
\[ -11213141 \div 18 \approx -622952.2778 \]
The integer part of the division is $-622953$.
\[ -11213141 = 18 \times (-622953) + r \]
where $r$ is the remainder.

2. Calculate the remainder $r$:
\[ r = -11213141 - 18 \times (-622953) \]
\[ r = -11213141 + 11213154 \]
\[ r = 3 \]

Therefore,
\[ -11213141 \equiv 3 \pmod{18} \]

So, the integer $n$ that satisfies $0 \le n < 18$ and $n \equiv -11213141 \pmod{18}$ is $3$.
\end{quote}

\paragraph{Turn 3 reply.} \texttt{Confidence: 95}

\paragraph{What happened.}
The method is right and every intermediate is right: $18 \times 622953 = 11213154$ checks out,
and the model has written the correct subtraction on the line above its answer.
It then evaluates $-11213141 + 11213154$ as $3$ rather than $13$. Gold is $13$.
Nothing in the presentation marks the slip. The derivation is well formed, correctly
structured, and confidently laid out, which is precisely why stated confidence rises
from $90$ to $95$ across an attempt that failed.

\paragraph{Readouts of this attempt.}
\begin{center}
\small
\begin{tabular}{lc}
\toprule
\textbf{readout} & \textbf{value} \\
\midrule
stated confidence, prospective (turn 1) & 90 \\
stated confidence, retrospective (turn 3) & 95 \\
P(True) & 0.32 \\
mean answer log-probability & $-0.05$ \\
\midrule
graded outcome & \textbf{incorrect} \\
\bottomrule
\end{tabular}
\end{center}

\noindent
This is one record, chosen to illustrate the failure mode rather than to carry evidential
weight; representative cell-level numbers are in \cref{tab:knowing}, \cref{app:cells}.

\section{Calibration beyond ECE, and cell-level multiplicity}
\label{app:stats}

ECE depends on binning and on sample composition, so we report two binning-free proper
scoring rules alongside it. On Mistral $\times$ MATH (750 records, accuracy 0.648), the
model's retrospective stated confidence gives Brier $0.262$ and NLL $1.470$ against a
10-bin ECE of $0.271$. P(True) on the same records gives Brier $0.386$: better ordered
than stated confidence, worse calibrated in absolute terms. The steering results in
\cref{tab:ablation} are reported as ECE because the steered runs stored per-gain
aggregates rather than per-record confidences; recomputing Brier and NLL at each gain
requires re-running the intervention and is the first thing we would add.

The cell-level bootstrap counts in \cref{tab:knowing} are uncorrected. With 16 cells,
some individually significant results are expected under the null, so the cell counts
should be read as descriptive. The aggregate pattern is the claim that carries weight:
the effect is positive in 12 of 16 cells with one tie, significantly positive in 10,
and never significantly negative. A
hierarchical model pooling across cells with model and domain as random effects is the
right test and is not yet run.

\section{Alternative explanations for the probe signal}
\label{app:alternatives}
A probe could key on properties of failed derivations rather than on a dedicated
metacognitive variable: low token probability, hedging language, or, on MATH, an
early-layer channel next to the answer text. Our surface-features baseline (answer
length, digit count, line count, word count, brace count) bounds this without closing
it, and the bound is domain dependent: adding the probe to the surface features raises
AUROC by $+0.21$ on MMLU-Pro but only $+0.07$ on MATH and $-0.11$ on geometry, where a
construction that fails to compile is visibly malformed. Two results argue against a pure surface account on the
domains where the signal is strong. A direction fitted on MMLU-Pro, where surface
reaches only 0.54 against a probe of 0.75, still governs MATH's self-report under
ablation ($0.880 \to 0.709$). And the pre-attempt read, taken where the activation is
byte-identical across a question's attempts, cannot recover the post-attempt direction
on MMLU-Pro (0.490, chance), which a surface-of-the-question account would predict it
could. A linear probe at one token also undercounts nonlinear or multi-token signal, so
the reported numbers are a floor.

\section{Per-cell numbers and specificity controls}
\label{app:cells}
\begin{table}[H]
\caption{Post-attempt correctness discrimination (AUROC). Left: self-report vs.\ the probe in representative cells, and the same comparison holding the question fixed (``within''), where the model's yardstick is P(True) across the five attempts at one question; overall the probe wins 12 of 16 with one tie, mean $+0.09$, with a paired bootstrap significantly positive in 10 of 16 and never significantly negative, uncorrected at the cell level (\cref{app:stats}). Right: Mistral specificity controls. ``pre'' is the last prompt token before the attempt; ``post'' is after the attempt; ``resid.'' removes the pre score; ``post$\to$pre'' applies that fitted post readout at the pre-attempt site, where only difficulty is visible.}
\label{tab:knowing}
\centering
\scriptsize
\begin{tabular}{@{}lcccc@{}}
\toprule
\textbf{model $\times$ domain} & \textbf{self-report} & \textbf{probe} & \textbf{within: P(True)} & \textbf{within: probe} \\
\midrule
Gemma-4 $\times$ MATH & 0.57 & \textbf{0.78} & -- & -- \\
Qwen3.6 $\times$ MMLU-Pro & 0.67 & \textbf{0.84} & 0.47 & \textbf{0.74} \\
GLM $\times$ MMLU-Pro & 0.67 & \textbf{0.85} & 0.59 & \textbf{0.73} \\
Mistral $\times$ MATH & 0.79 & \textbf{0.83} & 0.59 & \textbf{0.72} \\
\bottomrule
\end{tabular}
\hspace{0.9em}
\begin{tabular}{@{}lcccc@{}}
\toprule
\textbf{Mistral} & \textbf{pre} & \textbf{post} & \textbf{resid.} & \textbf{post$\to$pre} \\
\midrule
MMLU-Pro & 0.568 & 0.754 & 0.751 & 0.490 \\
MATH & 0.792 & 0.834 & 0.701 & 0.624 \\
GPQA & 0.593 & 0.572 & 0.544 & 0.479 \\
GeoGen & 0.770 & 0.676 & 0.505 & 0.648 \\
\\[-0.85em]
\bottomrule
\end{tabular}
\end{table}

\section{Checkpoints}
\label{app:checkpoints}
All models are public Hugging Face releases, run in bfloat16 without quantization, greedy decoding:
\texttt{mistralai/Mistral-Small-24B-Instruct-2501},
\texttt{Qwen/Qwen3.6-27B},
\texttt{zai-org/GLM-4.7-Flash},
\texttt{google/gemma-4-26B-A4B-it};
the dense within-family control for the lens is \texttt{Qwen/Qwen2.5-14B-Instruct}.

